% TEMPLATE-MILC-v3.tex - plantilla maestra UNICA de las guias MILC v3.
%
% La usan las cuatro familias de guias, cada una con su driver:
%   · Grados 6-11          -> scripts/build-guias-g11.py
%   · Bebras               -> scripts/build-guias-bebras.py
%   · Semillero            -> scripts/build-guias-semillero.py
%   · Territorio interior  -> scripts/build-guias-territorio-interior.py
%
% Antes habia tres copias casi identicas (~400 lineas iguales, 6 distintas):
% cada mejora habia que aplicarla tres veces, y en la practica no se hacia
% (el motor de recursos vivia solo en dos de ellas). Lo que varia entre
% familias esta parametrizado en 6 placeholders que cada driver llena
% (se nombran aqui SIN su sintaxis real: el builder reemplaza por texto plano,
% tambien dentro de los comentarios, y un valor multilinea rompe el preambulo):
%
%   HEADER_IZQ        encabezado izquierdo de pagina
%   HEADER_DER        encabezado derecho de pagina
%   PORTADA_ETIQUETA  rotulo sobre el numero grande (GRADO/MOMENTO/MODULO)
%   PORTADA_SERIE     linea de serie de la portada
%   PORTADA_ID        identificador de la guia en la portada
%   PORTADA_CIRCULO   el circulito de grado (°), o vacio si no aplica
%   PDF_TITULO        titulo en texto plano (sin \textbf ni comillas TeX) para
%                     los metadatos del PDF (/Title); lo produce lib_xelatex.texto_pdf
%
% Composicion (auditoria 2026-09-03): las bandas de estacion piden espacio
% (\Needspace*) y prohiben el salto justo despues (\nopagebreak), asi no quedan
% huerfanas al pie; las cajas largas (softbox, drawbox, infoband, puentebox) son
% `breakable`, asi una caja de media pagina ya no arrastra todo a la siguiente
% dejando la anterior al 40 %. Todo texto sobre mostaza o verde va en milcNegro
% (blanco sobre #E5B400 da 1,93:1 y sobre #58A923 2,95:1; ninguno pasa WCAG AA).
% El resto de claves las documenta content/guias/_SCHEMA.md. El script valida
% que no queden placeholders sin reemplazar antes de escribir el archivo final.

% Etiquetado PDF (latex-lab): /Lang, estructura y texto alternativo de las
% figuras, para lectores de pantalla. Debe ir ANTES de \documentclass.
\DocumentMetadata{lang=es-CO,tagging=on}
\documentclass[11pt,letterpaper]{article}
% headheight=24pt: la cabecera derecha lleva dos lineas (identidad de la guia
% y estacion actual). headsep se acorta en la misma medida para que el bloque
% de texto quede exactamente donde estaba con headheight=12pt.
\usepackage[margin=2.0cm,headheight=24pt,headsep=13pt]{geometry}
\usepackage[table]{xcolor}
\usepackage{fontspec}
\usepackage[spanish]{babel}
\usepackage{tikz}
% Librerías TikZ para los diagramas de las guías (bloque `recursos:`):
%  · babel  → babel en español vuelve activos caracteres como «>», y TikZ los usa
%             en sus opciones (>=stealth, ->). Sin esta librería, cualquier
%             diagrama con flechas revienta con «\language@active@arg> has an
%             extra }». Debe ir ANTES de que se dibuje nada.
%  · positioning        → sintaxis «below left=2pt and 3pt».
%  · decorations.pathreplacing → llaves ({brace}) para señalar rangos.
\usetikzlibrary{babel,positioning,decorations.pathreplacing}
\usepackage{graphicx}
% Circuitos: el trabajo de electrónica se hace hoy con micro:bit y sus sensores
% integrados (MakeCode, sin cableado), así que no se necesitan esquemáticos.
% Si algún día entran montajes con protoboard, basta descomentar esta línea:
% los diagramas .tex/.tikz declarados en `recursos:` ya se incrustan solos.
% \usepackage{circuitikz}
\usepackage[most]{tcolorbox}
\usepackage{tabularx}
\usepackage{array}
\usepackage{fancyhdr}
\usepackage{titlesec}
\usepackage[document]{ragged2e}  % bandera a la derecha en todo el cuerpo
\usepackage{enumitem}
\usepackage{microtype}
\usepackage{etoolbox}
\usepackage{needspace}
% hyperref al final del preambulo: metadatos del PDF (titulo, autor, idioma,
% familia), marcadores por estacion (\pdfbookmark) y enlaces sin recuadro.
\usepackage[hidelinks,
  pdftitle={Lo que controlo y lo que no: la pregunta que cambia la estrategia},
  pdfauthor={PhD. Álvaro Cárdenas Orozco},
  pdfsubject={Territorio interior · Educación socioemocional · Grado 10° · Momento 2 · MILC v3},
  pdflang={es-CO},
  bookmarksopen=true,bookmarksnumbered=false]{hyperref}

% Helvetica Neue no trae la flecha U+2192, asi que cada «→» escrito literalmente
% en el YAML se compilaba a NADA, en silencio (462 flechas en 61 guias: rutas del
% tipo «paleta → tipografia → lockup» salian rotas). Mapeamos el caracter a su
% version matematica, que si tiene glifo. Asi el autor puede seguir escribiendo
% «→» comodamente en el YAML.
\usepackage{newunicodechar}
\newunicodechar{→}{\ensuremath{\rightarrow}}
% Mismo problema con los simbolos de marca y vinetas: el «✓» de «una palomita ✓»
% salia como cuadrito vacio en el PDF de todas las guias que lo usaban.
\usepackage{pifont}
\newunicodechar{✓}{\ding{51}}
\newunicodechar{✗}{\ding{55}}
\newunicodechar{✔}{\ding{52}}
\newunicodechar{•}{\textbullet}
\newunicodechar{≥}{\ensuremath{\geq}}
\newunicodechar{≤}{\ensuremath{\leq}}

\setmainfont{Helvetica Neue}
\setsansfont{Helvetica Neue}
\setlength{\parindent}{0pt}
\setlength{\parskip}{6pt}
% Sin particion de palabras: en 6.º-7.º una palabra cortada por guion al final
% de linea («Comuni-cacion») frena la lectura mucho mas que una linea corta.
\hyphenpenalty=10000
\exhyphenpenalty=10000
\pretolerance=2000
\tolerance=3000
\setlength{\emergencystretch}{3em}
\linespread{1.10}
\renewcommand{\arraystretch}{1.25}
\setlist{leftmargin=5mm,itemsep=1mm,topsep=1mm}
\newcolumntype{Y}{>{\RaggedRight\arraybackslash}X}

\definecolor{milcMagenta}{HTML}{E6007E}
\definecolor{milcVino}{HTML}{5A0038}
\definecolor{milcTurquesa}{HTML}{0093A5}
\definecolor{milcVerde}{HTML}{58A923}
\definecolor{milcMostaza}{HTML}{E5B400}
\definecolor{milcOcre}{HTML}{B86600}
\definecolor{milcNegro}{HTML}{241816}
\definecolor{milcGris}{HTML}{F7F7F4}
\definecolor{milcLinea}{HTML}{E8E2DD}
\definecolor{milcGradePrimary}{HTML}{0D9488}
\definecolor{milcGradeDark}{HTML}{115E59}
\definecolor{milcGradeSoft}{HTML}{CCFBF1}
\definecolor{milcGradeAccent}{HTML}{CFE7FF}
\definecolor{milcGradeText}{HTML}{FFFFFF}
\color{milcNegro}

\pagestyle{fancy}
\fancyhf{}
% Cabecera de dos lineas: identidad de la guia arriba y, a la derecha y en
% gris, la estacion en curso (\markright desde \milcestacion). Antes de la
% primera estacion la segunda linea va vacia; el \strut mantiene la altura
% para que la linea de arriba no baile entre paginas.
\lhead{\small Territorio interior · Educación socioemocional\\[-1pt]{\footnotesize\strut}}
\rhead{\small Grado 10° · Momento 2 · MILC v3\\[-1pt]{\footnotesize\strut\color{milcNegro!65}\rightmark}}
\cfoot{\small\thepage}
\renewcommand{\headrulewidth}{0pt}

% Sobre mostaza (#E5B400) y verde (#58A923) el texto va en milcNegro; sobre el
% resto de la paleta, en blanco. Se usa dentro de fonttitle/fontupper, donde un
% \color posterior manda sobre coltitle.
\newcommand{\milctintasobre}[1]{%
  \ifboolexpr{test{\ifstrequal{#1}{milcMostaza}} or test{\ifstrequal{#1}{milcVerde}}}%
    {\color{milcNegro}}{\color{white}}}

\newtcolorbox{titlebox}[1]{
  enhanced,colback=#1,colframe=#1,arc=7mm,boxrule=0pt,
  left=7mm,right=7mm,top=3mm,bottom=3mm,
  before skip=8pt,after skip=8pt,
  fontupper=\sffamily\bfseries\Large\color{white}
}

\newtcolorbox{softbox}[3]{
  enhanced,breakable,pad at break=3mm,
  colback=#2,colframe=#1,borderline west={4pt}{0pt}{#1},
  arc=2mm,boxrule=.4pt,left=4mm,right=4mm,top=3mm,bottom=3mm,
  before skip=6pt,after skip=8pt,title={#3},coltitle=white,
  colbacktitle=#1,fonttitle=\sffamily\bfseries\milctintasobre{#1},fontupper=\RaggedRight,
  attach boxed title to top left={xshift=3mm,yshift=-2mm},
  boxed title style={arc=2mm, boxrule=0pt}
}

\newtcolorbox{drawbox}[2]{
  enhanced,breakable,pad at break=4mm,
  colback=white,colframe=#1,arc=4mm,boxrule=1pt,
  left=5mm,right=5mm,top=4mm,bottom=4mm,before skip=8pt,after skip=8pt,
  title={#2},coltitle=white,colbacktitle=#1,fonttitle=\sffamily\bfseries\milctintasobre{#1},
  fontupper=\RaggedRight,
  attach boxed title to top left={xshift=4mm,yshift=-2mm},
  boxed title style={arc=3mm, boxrule=0pt}
}

\newtcolorbox{infoband}[1]{
  enhanced,breakable,pad at break=2mm,colback=#1!8,colframe=#1!50,arc=2mm,boxrule=.5pt,
  left=4mm,right=4mm,top=2mm,bottom=2mm,before skip=4pt,after skip=4pt,
  fontupper=\small\RaggedRight
}

% Puente narrativo entre secciones (contrato editorial Conéctate).
% Da continuidad: "Acabas de X. Ahora vas a Y porque Z."
% Visualmente: banda izquierda en color del grado, texto en itálica pequeña.
\newtcolorbox{puentebox}{
  enhanced,breakable,pad at break=2mm,colback=milcGradeSoft,colframe=milcGradeDark,
  borderline west={3pt}{0pt}{milcGradeDark},
  arc=0mm,boxrule=0pt,left=5mm,right=4mm,top=2.5mm,bottom=2.5mm,
  before skip=4pt,after skip=8pt,
  fontupper=\itshape\small\color{milcNegro}\RaggedRight
}
% La flecha va en modo matematico a proposito: Helvetica Neue NO tiene el glifo
% U+2192, asi que \textrightarrow se compilaba a NADA («Missing character: There
% is no → in font Helvetica Neue») y el puente salia sin flecha. En modo
% matematico la flecha viene de la fuente matematica y si se dibuja.
\newcommand{\puente}[1]{%
  \begin{puentebox}%
    \textcolor{milcGradeDark}{$\rightarrow$}\hspace{2mm}#1%
  \end{puentebox}%
}

\newcommand{\linead}[1]{\noindent\color{milcLinea}\rule{#1}{0.5pt}\color{milcNegro}}
\newcommand{\respuesta}[1]{\par\vspace{2mm}\foreach \n in {1,...,#1}{\noindent\color{milcLinea}\rule{\linewidth}{0.5pt}\par\vspace{5mm}}\color{milcNegro}}
\newcommand{\checkbox}{\textcolor{milcGradeDark}{\fbox{\phantom{x}}}\hspace{5pt}}

% ─── Listas aireadas para las guías socioemocionales ────────────────────
% Componer las verificaciones como «\checkbox texto\\» deja el texto que salta
% de línea debajo del cuadrito y apelmaza el bloque. Estos entornos alinean la
% continuación y airean los ítems. Los usa Territorio Interior; cualquier
% familia puede hacerlo.
\newlist{checklista}{itemize}{1}
\setlist[checklista]{
  label=\textcolor{milcGradeDark}{\fbox{\phantom{x}}},
  leftmargin=9mm, labelsep=3.2mm, itemsep=2.4mm,
  topsep=2.5mm, parsep=0pt, partopsep=0pt, before=\RaggedRight
}
\newlist{pasoslista}{enumerate}{1}
\setlist[pasoslista]{
  label=\textcolor{milcGradeDark}{\sffamily\bfseries\arabic*.},
  leftmargin=8mm, labelsep=3mm, itemsep=2.8mm,
  topsep=1mm, parsep=0pt, partopsep=0pt, before=\RaggedRight
}

\newcommand{\phasecard}[4]{
\begin{tcolorbox}[enhanced,colback=#3,colframe=#2,arc=3mm,boxrule=.5pt,height=3.05cm,valign=center,halign=center,left=1.5mm,right=1.5mm,top=2mm,bottom=2mm]
{\sffamily\bfseries\fontsize{22}{24}\selectfont\textcolor{#2}{#1}}\par
{\sffamily\bfseries\small #4}
\end{tcolorbox}
}

% ── Piezas del rediseno 2026 (legibilidad 6.º-11.º) ───────────────────
% Banda de estacion: reemplaza el titlebox macizo. Titulo a la izquierda,
% dato practico (tiempo/modalidad) a la derecha, en una sola linea.
% Nunca huerfana: pide 9 lineas libres (o salta de pagina rellenando con
% \vfill) y prohibe el salto justo despues de la banda. Nueve y no seis:
% la caja partible que sigue necesita ~80pt (titulo adosado, relleno y dos
% lineas) para arrancar en la misma pagina; con menos, tcolorbox la manda
% entera a la siguiente y la banda se queda sola. El penalty 10000 va
% en `after`, ANTES del espacio posterior: un `after skip` (aunque sea 0pt)
% es un \vskip tras la caja, y ese glue es un punto de corte legal que un
% \nopagebreak escrito despues de \end{tcolorbox} ya no cierra. Cada banda
% deja marcador PDF y actualiza la cabecera derecha.
\newcounter{milcestacion}
\newcommand{\milcestacion}[2]{%
\Needspace*{9\baselineskip}%
\stepcounter{milcestacion}%
\pdfbookmark[1]{#2}{milcestacion\themilcestacion}%
\markright{#2}%
\begin{tcolorbox}[enhanced,colback=#1,colframe=#1,arc=2mm,boxrule=0pt,
  left=5mm,right=5mm,top=2.6mm,bottom=2.6mm,before skip=11pt,
  after={\par\nopagebreak\vskip7pt}]
{\sffamily\bfseries\fontsize{15}{18}\selectfont\milctintasobre{#1}#2}
\end{tcolorbox}}

% Tarjeta de paso numerada: sustituye las tablas «Pilar / Aplicacion», que
% eran formato de consulta docente, no de estudiante. El numero va en un
% circulo sobre el borde para que la secuencia se lea de un vistazo.
\newcommand{\milcpaso}[3]{%
\Needspace{4\baselineskip}%
\begin{tcolorbox}[enhanced,colback=white,colframe=milcLinea,arc=1.5mm,boxrule=.9pt,
  left=14mm,right=4mm,top=3mm,bottom=3mm,before skip=4pt,after skip=4pt,
  fontupper=\RaggedRight,
  overlay={\node[anchor=north west,circle,fill=#2,text=white,inner sep=0pt,
    minimum size=7.4mm,font=\sffamily\bfseries\small]
    at ([xshift=3.6mm,yshift=-3.2mm]frame.north west) {#1};}]
#3
\end{tcolorbox}}

% Casilla marcable de tamano util con lapiz (la anterior era \fbox{\phantom{x}},
% de alto variable segun el texto que la seguia).
\newcommand{\milcmarca}{\framebox[3.8mm]{\rule{0pt}{3.4mm}}}

% Fila marcable de lista suelta (checks de la estacion 1).
\newcommand{\milccheckrow}[1]{%
\par\vspace{1.8mm}\noindent
\makebox[6.4mm][l]{\raisebox{-.55mm}{\textcolor{milcNegro}{\milcmarca}}}%
\parbox[t]{\dimexpr\linewidth-6.4mm\relax}{\RaggedRight #1}\par}

% Celda marcable dentro de la rubrica (tabla de autoevaluacion). Misma
% sangria francesa, pero pensada para vivir dentro de una columna X.
\newcommand{\milcopcion}[1]{%
\makebox[5.2mm][l]{\raisebox{-.55mm}{\textcolor{milcNegro}{\milcmarca}}}%
\parbox[t]{\dimexpr\linewidth-5.2mm\relax}{\RaggedRight #1}}

% ── Recursos visuales de la guía ──────────────────────────────────────
% Declarados en el bloque `recursos:` del YAML y emitidos por el builder.
% \guiaFigura   → imagen rasterizada o PDF (foto, carta celeste, montaje).
% \guiaDiagrama → fuente TikZ/circuitikz (.tex) incrustada como vector:
%                 es la vía para circuitos y esquemáticos editables.
% [#1] = texto alternativo (alt) · #2 = ruta relativa al .tex · #3 = pie
% El alt llega al PDF etiquetado (/Alt) y lo lee el lector de pantalla.
\newcommand{\guiaFigura}[3][]{%
\begin{center}
\includegraphics[width=0.88\linewidth,height=0.38\textheight,keepaspectratio,alt={#1}]{#2}\\[1.5mm]
{\footnotesize\itshape\textcolor{milcNegro!75}{#3}}
\end{center}
}

\newcommand{\guiaDiagrama}[2]{%
\begin{center}
\input{#1}\\[1.5mm]
{\footnotesize\itshape\textcolor{milcNegro!75}{#2}}
\end{center}
}

\begin{document}
\thispagestyle{empty}

% ── Portada ───────────────────────────────────────────────────────────
% Antes: un rectangulo de color a pagina completa. Se veia bien en pantalla,
% pero en la fotocopiadora del colegio se comia el toner y salia gris sucio.
% Ahora la banda ocupa el tercio superior y el resto va en blanco.
\begin{tikzpicture}[remember picture,overlay]
  \path[fill=milcGradePrimary]
    (current page.north west) rectangle ([yshift=-10.6cm]current page.north east);

  \node[anchor=north west,text=milcGradeText] at ([xshift=2.0cm,yshift=-1.55cm]current page.north west)
    {\sffamily\bfseries\fontsize{12}{14}\selectfont MOMENTO};
  \node[anchor=north west,text=milcGradeText,opacity=.85] at ([xshift=2.0cm,yshift=-2.35cm]current page.north west)
    {\sffamily\bfseries\fontsize{8.5}{10}\selectfont TERRITORIO INTERIOR · SOCIOEMOCIONAL};
  \node[anchor=north east,text=milcGradeText,opacity=.85,align=right] at ([xshift=-2.0cm,yshift=-1.55cm]current page.north east)
    {\sffamily\bfseries\fontsize{9}{12}\selectfont I.E. Sor María Juliana\\Cartago, Valle del Cauca};

  \node[anchor=west,text=milcGradeText] (gradonum) at ([xshift=1.85cm,yshift=-7.1cm]current page.north west)
    {\sffamily\bfseries\fontsize{112}{112}\selectfont 2};
  % El circulito de grado (°) solo aplica a las guias de grado («11°»); en
  % Bebras y Semillero el numero grande es un momento/modulo, no un grado.
  % Se posiciona relativo al nodo `gradonum`, no en coordenadas absolutas.
  

  \node[anchor=west,text=milcGradeText,text width=11.4cm,align=left] at ([xshift=1.5cm,yshift=0.5cm]gradonum.east)
    {
     {\sffamily\bfseries\fontsize{9}{11}\selectfont\MakeUppercase{Territorio interior · Socioemocional}}\\[3mm]
     {\sffamily\bfseries\fontsize{21}{25}\selectfont\setlength{\baselineskip}{27pt}Lo que controlo\\[4pt]y lo que no\\[4pt]la pregunta que cambia todo}\\[3.5mm]
     {\sffamily\bfseries\fontsize{15}{18}\selectfont Grado 10° · Momento 2}
    };
\end{tikzpicture}

\vspace*{9.4cm}
\vfill

{\sffamily\bfseries\fontsize{8}{10}\selectfont\textcolor{milcGradeDark}{LO QUE TE LLEVAS HOY}}

\begin{tcolorbox}[enhanced,colback=white,colframe=milcNegro,arc=2mm,boxrule=1.6pt,
  left=5mm,right=5mm,top=4mm,bottom=4mm,before skip=3pt,after skip=12pt,
  fontupper=\RaggedRight\fontsize{11}{15.5}\selectfont]
Tu tabla de dos columnas: las diez cosas que te quitan el sueño repartidas entre lo que depende de ti y lo que no, con una acción concreta para cada una de la primera columna y una forma de sostener para cada una de la segunda, más la que estabas tratando con la estrategia equivocada.
\end{tcolorbox}

\begin{tcolorbox}[enhanced,colback=milcGris,colframe=milcGris,arc=2mm,boxrule=0pt,
  left=5mm,right=5mm,top=3.5mm,bottom=3.5mm,after skip=12pt]
{\sffamily\bfseries\fontsize{8}{10}\selectfont\textcolor{milcNegro!55}{ESTA GUÍA ES DE}}\\[3mm]
\begin{tabularx}{\linewidth}{X p{3.2cm} p{3.2cm}}
\linead{\linewidth} & \linead{3.0cm} & \linead{3.0cm}\\[-1mm]
{\fontsize{8.5}{10}\selectfont\textcolor{milcNegro!55}{Nombre completo}} &
{\fontsize{8.5}{10}\selectfont\textcolor{milcNegro!55}{Grupo}} &
{\fontsize{8.5}{10}\selectfont\textcolor{milcNegro!55}{Fecha}}\\
\end{tabularx}
\end{tcolorbox}

\vfill

\noindent\textcolor{milcLinea}{\rule{\linewidth}{1pt}}\\[2mm]
{\sffamily\bfseries\fontsize{9.5}{12}\selectfont Autor: Dr. Álvaro Cárdenas Orozco}\\[1mm]
{\sffamily\fontsize{8.5}{11}\selectfont\textcolor{milcNegro!55}{Metodología MILC · Apertura ancestral · Dicotomía del control · Triángulo Dussel-Estoicismo-Floridi}}

\newpage

\section*{Apertura: Lo que controlo y lo que no: la pregunta que cambia la estrategia}

\begin{softbox}{milcGradeDark}{milcGradeSoft}{Saber ancestral}
En \textbf{1983}, en Chinchiná (Caldas), se detectó por primera vez en Colombia la \textbf{roya del cafeto} ---un hongo que ataca la hoja y puede arrasar un cafetal---. \emph{Para las familias cafeteras de esta región no había nada que negociar con eso:} \textbf{el hongo llegó, no lo trajo nadie, y no se iba a ir porque alguien se preocupara}. \textbf{Y sin embargo la historia no termina ahí, y por eso es el ancla de esta guía.} En \textbf{1984} el Centro Nacional de Investigaciones de Café ---\textbf{Cenicafé}--- entregó la \textbf{Variedad Colombia}, resistente a la roya. \emph{Fíjate en el dato que lo vuelve extraordinario:} \textbf{esa variedad era el resultado de \emph{veinte años} de mejoramiento genético}, un trabajo empezado \textbf{cuando la roya todavía no había llegado al país}. \emph{Alguien decidió, dos décadas antes, prepararse para algo que aún no pasaba.} \textbf{Y quedó una decisión más, que sí era de cada caficultor: \emph{renovar o no renovar} el cafetal con materiales resistentes} ---renovar costaba, porque significaba perder cosechas mientras el cafetal nuevo crecía---. \emph{Ahí está el reparto entero: el hongo no dependía de nadie; la preparación y la decisión de renovar, sí.}
\end{softbox}

\begin{softbox}{milcTurquesa}{milcGris}{Saber contemporáneo}
¿Y sirve de algo saber si algo depende de uno? \textbf{Sirve para escoger la herramienta, que es de lo que depende que funcione.} \textbf{Richard Lazarus} y \textbf{Susan Folkman} describieron dos formas de afrontar lo que estresa. El \textbf{afrontamiento centrado en el problema} intenta \emph{cambiar la situación} que causa el malestar: planear, actuar, buscar información, pedir ayuda. El \textbf{afrontamiento centrado en la emoción} intenta \emph{manejar la respuesta emocional} ante algo que no se puede cambiar: aceptar, reinterpretar, buscar apoyo, distraerse a propósito (Lazarus y Folkman, 1984). \textbf{Y el hallazgo que hay que memorizar no es que una sea mejor que la otra ---es que cada una sirve para un tipo de situación---:} \emph{el afrontamiento centrado en el problema funciona cuando la situación se evalúa como \textbf{controlable o modificable}; el centrado en la emoción, cuando se evalúa como \textbf{inmodificable}}. \textbf{La eficacia de una estrategia depende, sobre todo, de qué tan controlable sea aquello a lo que se aplica.} \emph{Traducido a tu vida:} \textbf{el problema no suele ser que te esfuerces poco ---es que estás usando la herramienta equivocada---}, y eso agota sin producir nada.
\end{softbox}

\begin{softbox}{milcMostaza}{milcGris}{Pregunta-puente}
\textit{Cenicafé empezó a prepararse \textbf{veinte años antes} de que llegara la roya, y cada familia tuvo que decidir si renovaba. \textbf{De las cosas que te quitan el sueño, ¿cuál llevas meses tratando de controlar} \ldots \textbf{sabiendo que no depende de ti}?}
\end{softbox}

\begin{softbox}{milcVerde}{milcGris}{Saber y saber hacer}
Al terminar podrás: (1) \textbf{separar} lo que depende de ti de lo que no, con un criterio y no con una intuición; (2) \textbf{explicar} las dos formas de afrontamiento y cuándo sirve cada una; (3) \textbf{detectar} en qué preocupación estás usando la estrategia equivocada, que es de donde viene buena parte del agotamiento; (4) \textbf{escribir} una acción concreta para lo controlable y una forma de sostener para lo que no lo es.
\end{softbox}



\small
\begin{tabularx}{\textwidth}{>{\bfseries}p{3.0cm}Y}
\rowcolor{milcGradePrimary!12}Autor & Dr. Álvaro Cárdenas Orozco\\
DBA & Comprende los fenómenos que afectan a su territorio, regula sus emociones ante ellos y acompaña a otros con empatía (transversal 6°--11°, articulado con la política «Escuela, territorio de vida» del MEN, Resolución 006519 de 2025, y la Ley 1620 de 2013)\\
Referentes & Primera ayuda psicológica (OMS, 2011/2012) · Directrices IASC (2007) · USGS y Servicio Geológico Colombiano · Pueblo nasa y la minga (CRIC) · Dussel · Estoicismo · Floridi\\
Producto final & Tu tabla de dos columnas: las diez cosas que te quitan el sueño repartidas entre lo que depende de ti y lo que no, con una acción concreta para cada una de la primera columna y una forma de sostener para cada una de la segunda, más la que estabas tratando con la estrategia equivocada.\\
\end{tabularx}
\normalsize

\puente{En el momento 1 aprendiste a nombrar la ansiedad. \textbf{Hoy le pones método:} casi todo el desgaste viene de empujar lo que no se mueve.}

\milcestacion{milcGradeDark}{Ruta MILC: así vamos a aprender}
\begin{tabularx}{\textwidth}{>{\centering\arraybackslash}X>{\centering\arraybackslash}X>{\centering\arraybackslash}X>{\centering\arraybackslash}X}
\phasecard{1}{milcVerde}{milcGris}{Escucha\\[-1mm]\small Observo, escucho y pregunto.} &
\phasecard{2}{milcTurquesa}{milcGris}{Sistematización\\[-1mm]\small Separo y actúo donde puedo.} &
\phasecard{3}{milcMagenta}{milcGris}{Praxis\\[-1mm]\small Construyo y aplico.} &
\phasecard{4}{milcVino}{milcGris}{Evaluación\\[-1mm]\small Reflexión liberadora.}
\end{tabularx}


\puente{Primero, la lista honesta de lo que te quita el sueño. \emph{Sin ordenarla todavía.}}

\milcestacion{milcVerde}{Estación 1 de 3 · Escucha}

\begin{softbox}{milcVerde}{milcGris}{Escena para escuchar}
\textbf{Actividad 1 · IDENTIFICA --- Lo que me quita el sueño.} \textbf{Nadie va a leer esta página.} \textbf{Primera parte.} Escribe \textbf{diez cosas} que te preocupan de verdad. \emph{De todos los tamaños:} el examen del viernes, qué vas a estudiar, la plata en la casa, si alguien está enfermo, cómo te ves, si te van a dejar de hablar, el país, el futuro. \textbf{Segunda parte.} Al frente de cada una, escribe \textbf{qué has hecho} con esa preocupación esta semana. \emph{Y sé literal:} «pensar en eso antes de dormir» \textbf{es una respuesta válida y es el dato}. \textbf{Tercera parte.} Marca en cuáles \textbf{hiciste algo} y en cuáles \textbf{solo le diste vueltas}. \textbf{Y la última:} de las que solo diste vueltas, \emph{¿había algo que pudieras hacer?} \emph{Contesta con honestidad: a veces sí y no lo hiciste, y a veces de verdad no había nada. Son casos muy distintos.}
\end{softbox}

{\sffamily\bfseries\fontsize{8}{10}\selectfont\textcolor{milcNegro!55}{MARCA CUANDO LO HAYAS HECHO}}
\milccheckrow{Escribiste diez preocupaciones reales, de todos los tamaños.}
\milccheckrow{Anotaste \textbf{qué hiciste} con cada una esta semana, aunque fuera darle vueltas.}
\milccheckrow{Separaste en cuáles actuaste y en cuáles solo pensaste.}

\begin{infoband}{milcVerde}


\textbf{En tu cuaderno (Actividad 1 · IDENTIFICA).} Título: «Actividad 1. Lo que me quita el sueño». \textbf{Formato:} las diez preocupaciones + qué hiciste con cada una + la marca de acción o vueltas. \textbf{Extensión:} media página. \textbf{Sabes que terminaste cuando} puedes decir \textbf{cuántas de las diez son de darle vueltas}. Esta página es tuya: \textbf{nadie más tiene que leerla si tú no quieres}.
\end{infoband}

\puente{Ya la tienes. Ahora, las dos estrategias, cuándo sirve cada una y por qué usar la equivocada agota tanto.}

\milcestacion{milcTurquesa}{Fase 2 · Sistematización: entender el control}

\begin{softbox}{milcTurquesa}{milcGris}{Cuatro claves sobre lo que depende de ti}
Cuatro claves para dejar de gastar fuerza donde no rinde.
\end{softbox}

\milcpaso{1}{milcTurquesa}{\textbf{Hay dos estrategias, y las dos son válidas.} \textbf{Centrada en el problema:} cambiar la situación ---planear, actuar, informarse, pedir ayuda---. \textbf{Centrada en la emoción:} manejar la respuesta ante algo que no cambia ---aceptar, reinterpretar, buscar apoyo, distraerse a propósito--- (Lazarus y Folkman, 1984). \emph{Mucha gente cree que la segunda es «rendirse», y no lo es:} \textbf{cuando algo de verdad no depende de ti, cuidar cómo lo llevas \emph{es} la acción correcta}. \emph{Y al revés:} \textbf{cuando algo sí depende de ti, calmarte y no hacer nada es evitación con buen nombre.}}
\milcpaso{2}{milcTurquesa}{\textbf{Usar la equivocada agota, y esa es la causa del cansancio raro.} La eficacia depende de \textbf{qué tan controlable} sea la situación: lo centrado en el problema sirve para lo \textbf{modificable}; lo centrado en la emoción, para lo \textbf{inmodificable}. \emph{Junta eso con tu lista y vas a ver el patrón:} \textbf{la preocupación que más te desgasta suele ser una en la que llevas meses aplicando la herramienta equivocada.} \emph{Los dos errores típicos:} \textbf{planear y planear sobre algo que no depende de ti} ---la salud de alguien, lo que otro piense, lo que pase en el país--- \emph{y} \textbf{«aceptar» algo que sí podías cambiar} ---una materia que vas perdiendo, una conversación que no has tenido---.}
\milcpaso{3}{milcTurquesa}{\textbf{Casi nada es todo o nada: hay que partir la preocupación.} \emph{La mayoría de las cosas grandes tienen \textbf{una parte} que depende de ti, y por eso el reparto no se hace con la preocupación entera sino con sus pedazos.} \textbf{Ejemplo:} \emph{«la plata en mi casa» no depende de ti ---pero \textbf{gastar menos}, \textbf{no pedir cosas caras} o \textbf{ayudar con algo} sí---.} \textbf{Otro:} \emph{«si me va a dar la nota» no depende solo de ti ---pero \textbf{cuántas horas estudias} y \textbf{si preguntas lo que no entendiste} sí---.} \emph{La pregunta útil no es «¿esto depende de mí?» sino \textbf{«¿qué parte de esto depende de mí?»}.} \textbf{Casi siempre hay una, y casi siempre es más pequeña de lo que uno querría y más grande de lo que uno cree.}}
\milcpaso{4}{milcTurquesa}{\textbf{Prepararse antes es la jugada más poderosa, y casi nadie la usa.} \emph{Cenicafé empezó a buscar una variedad resistente \textbf{veinte años antes} de que la roya llegara al país.} \textbf{Cuando llegó, el país no estaba improvisando: tenía una respuesta lista un año después.} \emph{Eso tiene una traducción muy concreta a tu vida:} \textbf{hay cosas que hoy no dependen de ti \emph{porque ya pasaron}, pero cuya próxima vez sí depende de lo que hagas ahora} ---estudiar antes en vez de la noche anterior, aprender algo antes de necesitarlo, tener una conversación difícil antes de que explote, saber a quién acudir antes de la crisis---. \emph{Prepararse es convertir, con tiempo, algo incontrolable en algo controlable.}}

\begin{drawbox}{milcTurquesa}{El reparto, en cuatro preguntas}
\begin{pasoslista}
\item \textbf{¿Qué parte de esto depende de mí?} \emph{No la preocupación entera: sus pedazos.}
\item \textbf{Si depende:} \emph{¿cuál es la siguiente acción concreta, y cuándo la hago?}
\item \textbf{Si no depende:} \emph{¿cómo lo sostengo?} ---contárselo a alguien, moverse, dormir, no rumiarlo antes de dormir---.
\item \textbf{¿Estoy usando la herramienta correcta} para esta columna?
\item \textbf{Y la de largo plazo:} \emph{¿hay algo que pueda hacer hoy para que la próxima vez sí dependa de mí?}
\end{pasoslista}
\end{drawbox}

\begin{softbox}{milcMostaza}{milcGris}{Errores comunes y su corrección}
\textbf{Creer que preocuparse es hacer algo:} darle vueltas no cambia nada y cansa igual. \textbf{Repartir la preocupación entera:} casi todas tienen una parte tuya. \textbf{Llamar aceptación a la evitación:} si podías actuar, no era aceptar. \textbf{Planear sobre lo que no depende de ti:} es el desgaste más común. \textbf{Creer que lo emocional es rendirse:} ante lo inmodificable, es la estrategia correcta. \textbf{No preparar la próxima vez:} lo incontrolable de hoy suele venir de algo que no se preparó. \textbf{Hacer la lista y no escribir acciones:} una tabla sin acciones es la misma preocupación en columnas.
\end{softbox}

\begin{infoband}{milcTurquesa}


\textbf{En tu cuaderno (Actividad 2 · EXPLICA).} Título: «Actividad 2. Las dos estrategias». \textbf{Formato:} define las dos con tus palabras y pon un ejemplo tuyo de cada una \textbf{bien usada} y uno de cada una \textbf{mal usada}. \textbf{Extensión:} media página. \textbf{Sabes que terminaste cuando} puedes explicar por qué \textbf{la eficacia depende de qué tan controlable sea la situación}.
\end{infoband}

\puente{Ya sabes el criterio. Es hora de repartir en dos columnas y darle a cada una lo que le sirve.}

\milcestacion{milcMagenta}{Fase 3 · Praxis: repartir en dos columnas}

\begin{softbox}{milcMagenta}{milcGris}{Cómo se reparte lo que quita el sueño}
Una tabla de dos columnas, con algo escrito al frente de cada línea. Ese «algo» es lo que la vuelve distinta de una lista de preocupaciones.
\end{softbox}

\milcpaso{1}{milcMagenta}{\textbf{Las diez preocupaciones (5 min).} Pásalas en limpio. \emph{Y párte las que sean muy grandes:} \textbf{«el futuro» no se puede repartir; «qué voy a estudiar el otro año» sí}.}
\milcpaso{2}{milcMagenta}{\textbf{El reparto (10 min).} Dos columnas: \textbf{depende de mí} · \textbf{no depende de mí}. \emph{Y aplica la regla del pilar 3:} \textbf{si algo se te va entero a la derecha, búscale el pedazo que sí es tuyo y súbelo a la izquierda}. \textbf{Casi todas las líneas terminan partidas en dos, y eso está bien: es lo que se busca.}}
\milcpaso{3}{milcMagenta}{\textbf{Una acción y una forma de sostener (12 min).} \emph{Columna izquierda:} al frente de cada línea, \textbf{la siguiente acción concreta con día} ---no «estudiar más», sino «el martes le pregunto al profe las dos cosas que no entendí»---. \emph{Columna derecha:} al frente de cada línea, \textbf{cómo lo vas a sostener} ---a quién se lo cuentas, qué haces cuando aparece de noche, qué te ayuda de verdad---. \textbf{Y una regla para la derecha:} \emph{«no pensar en eso» no es una forma de sostener; es evitación, y ya sabes lo que hace}.}
\milcpaso{4}{milcMagenta}{\textbf{La estrategia equivocada (8 min).} Señala \textbf{la preocupación en la que llevas más tiempo usando la herramienta que no era} y escribe \textbf{qué vas a hacer distinto}. \emph{Y agrega la jugada de Cenicafé:} \textbf{escoge una cosa que hoy no dependa de ti y escribe qué podrías empezar a preparar ahora para que la próxima vez sí dependa}. \emph{Con fecha, aunque el plazo sea largo.}}

\begin{drawbox}{milcMagenta}{Antes de cerrar tu tabla}
\begin{checklista}
\item Partiste las preocupaciones grandes en pedazos manejables.
\item Cada línea de la izquierda tiene \textbf{una acción concreta con día}.
\item Cada línea de la derecha tiene \textbf{una forma de sostener} ---y no es «no pensar en eso»---.
\item Buscaste el pedazo propio de las que se iban enteras a la derecha.
\item Señalaste la preocupación en la que usabas la herramienta equivocada.
\item Escribiste \textbf{una cosa que vas a preparar} para que la próxima vez sí dependa de ti.
\end{checklista}
\end{drawbox}

\begin{softbox}{milcMagenta}{milcGris}{Plantilla de trabajo}
\textbf{El reparto.} \emph{«De \_\_\_\_, lo que depende de mí es \_\_\_\_ y lo que no es \_\_\_\_.»} \\[1mm]
\textbf{Izquierda.} \emph{«Siguiente acción: \_\_\_\_. Día: \_\_\_\_.»} \\[1mm]
\textbf{Derecha.} \emph{«Lo sostengo \_\_\_\_. Se lo cuento a \_\_\_\_.»} ---nunca «no pensar en eso»---.
\end{softbox}

\begin{infoband}{milcMagenta}


\textbf{En tu cuaderno (Actividad 3 · CREA).} Título: «Actividad 3. Mi tabla de dos columnas». \textbf{Formato:} la tabla completa con acción o forma de sostener al frente de cada línea + la preocupación con la estrategia equivocada + lo que vas a preparar. \textbf{Extensión:} una página. \textbf{Sabes que terminaste cuando} \textbf{ninguna línea quedó sin nada escrito al frente}.
\end{infoband}

\puente{El producto es una tabla con acciones. \emph{Una preocupación sin acción al lado sigue siendo una preocupación.}}

\milcestacion{milcMostaza}{Producto de la sesión}
\textbf{Tabla de dos columnas: lo que depende de ti con su acción, lo que no con su forma de sostener}.
\begin{softbox}{milcMostaza}{milcGris}{Criterios de calidad}
(1) Las preocupaciones grandes están partidas en pedazos y no se reparten enteras. (2) Cada línea de la columna controlable tiene una acción concreta y fechada. (3) Cada línea de la columna no controlable tiene una forma de sostener que no sea evitación. (4) Se identifica al menos una preocupación tratada con la estrategia equivocada, y qué se hará distinto. (5) Se planea algo que hoy no depende de la persona y que puede empezar a prepararse.
\end{softbox}

\puente{Antes de cerrar, revisa lo que casi todos hacen mal: si en la segunda columna escribiste acciones para cambiar la situación, sigues empujando lo que no se mueve.}

\milcestacion{milcVino}{Evaluación liberadora · cinco dimensiones}

\begin{softbox}{milcVino}{milcGris}{Sentido de la evaluación}
Evaluar no es solo revisar una nota. Es reconocer qué aprendí, qué cambió en mí, qué emociones manejé, cómo me relaciono con la ciudad y la comunidad, y qué le dejaré a quien viene detrás.
\end{softbox}

\textbf{Mi reflexión liberadora en cinco dimensiones.} Completa cada frase iniciada:

\small
\begin{tabularx}{\textwidth}{>{\bfseries\RaggedRight\arraybackslash}p{3.4cm}Y>{\RaggedRight\arraybackslash}p{4.6cm}}
\rowcolor{milcVino}\color{white}Dimensión & \color{white}\bfseries Mi respuesta (frame starter) & \color{white}\bfseries Algo que descubrí\\
1. Personal & Lo que más cambió en mí fue \linead{6.5cm} & \linead{4.2cm}\\
2. Emocional & La emoción más fuerte fue \linead{2.5cm} y la regulé \linead{3cm} & \linead{4.2cm}\\
3. Ciudadana & En democracia esto importa porque \linead{6cm} & \linead{4.2cm}\\
4. Local (Cartago) & En mi barrio sería útil para \linead{6.5cm} & \linead{4.2cm}\\
5. Intergeneracional & A mi abuela/abuelo le diría \linead{6.5cm} & \linead{4.2cm}\\
\end{tabularx}
\normalsize

\begin{infoband}{milcVino}
\textbf{Para la próxima sustentación.} Vuelve a esta tabla en seis meses. Verás que las respuestas cambian — no porque lo aprendido cambie, sino porque tú cambias. La evaluación liberadora es una conversación contigo a través del tiempo.
\end{infoband}


\puente{Cerramos con tres voces: el que está más allá de la totalidad, lo que sí está en tu poder, y los pequeños patrones que arman un juicio.}

\milcestacion{milcGradeDark}{Triángulo de pensamiento · cierre filosófico}

\begin{softbox}{milcVino}{milcGris}{Dussel · lente del nosotros}
\textit{«Analéctico quiere indicar el hecho real humano por el que todo hombre, todo grupo o pueblo se sitúa siempre más allá (aná-) del horizonte de la totalidad.»}

{\footnotesize\itshape\color{milcNegro!70}--- Enrique Dussel · Filosofía de la liberación (1977), §2.4}

Hay una lectura de «lo que no depende de mí» que es cómoda y falsa: \emph{usarla para no hacer nada nunca}. \textbf{Dussel ayuda a evitarla, porque su idea es justamente que nadie está encerrado del todo en la situación que le tocó:} \emph{siempre se está \textbf{más allá} del horizonte que lo rodea}. \textbf{Aplícalo a la columna derecha:} \emph{muchas cosas que se sienten como destino ---la plata que hay en la casa, el barrio, lo que se espera de ti--- \textbf{son historia, no naturaleza}}, y la historia la hace gente. \textbf{Cenicafé es la prueba pequeña:} \emph{la roya era inevitable; que llegara y no hubiera respuesta, no lo era.}

\textbf{De lo que puse en la columna derecha, ¿qué es de verdad inevitable y qué solo lleva mucho tiempo siendo así?}
\end{softbox}
\linead{\linewidth}\\[1mm]
\linead{\linewidth}

\begin{softbox}{milcMostaza}{milcGris}{Estoicismo · Epicteto · Enquiridión, 1 (c. 125 d.C.; trad. desde E. Carter) · cuidado interior}
\textit{«Algunas cosas están en nuestro poder y otras no. Están en nuestro poder la opinión, el impulso, el deseo, el rechazo; en una palabra, todo aquello que es obra nuestra.»}

{\footnotesize\itshape\color{milcNegro!70}--- Epicteto · Enquiridión, 1 (c. 125 d.C.; trad. desde E. Carter)}

Esta es la frase que funda toda la guía, y conviene leerla con precisión. \emph{Epicteto no dice que lo de afuera no importe ---dice que no es \textbf{obra nuestra}---.} \textbf{Y fíjate en lo que sí pone del lado nuestro: la opinión, el impulso, el deseo, el rechazo.} \emph{Es decir: lo que pienso de lo que pasa y lo que hago con ello.} \textbf{Ahora la advertencia, porque esta frase se usa mal muy seguido:} \emph{no sirve para decirle a alguien que su problema no es real}. \textbf{Que el hambre, la enfermedad o la violencia no dependan de una persona no las vuelve menos graves ---las vuelve \emph{injustas}, que es otra cosa---.} \emph{La dicotomía sirve para escoger dónde poner la fuerza, no para negar lo que duele.}

\textbf{¿Estoy usando «no depende de mí» para escoger bien, o para no hacer lo que sí puedo?}
\end{softbox}
\linead{\linewidth}\\[1mm]
\linead{\linewidth}

\begin{softbox}{milcTurquesa}{milcGris}{Floridi · lente de la infoesfera}
\textit{«Los pequeños patrones solo pueden ser significativos si se agregan correctamente, se comparan y se procesan a tiempo.»}

{\footnotesize\itshape\color{milcNegro!70}--- Luciano Floridi · Big data and their epistemological challenge (2012)}

Hay una forma moderna de gastar fuerza en la columna derecha y tiene nombre: \textbf{informarse compulsivamente}. \emph{Leer todo el día sobre algo que no puedes cambiar ---el país, una guerra, una enfermedad--- se siente como estar haciendo algo}, \textbf{y no lo es}. \emph{Floridi señala por qué además no informa:} \textbf{los patrones pequeños solo significan algo si se agregan bien y \emph{se procesan a tiempo}}; \emph{un chorro continuo de titulares no se procesa ---se acumula---}. \textbf{La versión práctica:} \emph{para lo de la columna derecha, informarse una vez al día en una fuente buena \textbf{informa más} que revisar veinte veces}, y deja tiempo para lo de la izquierda.

\textbf{¿Cuánto tiempo al día le dedico a leer sobre cosas que no dependen de mí?}
\end{softbox}
\linead{\linewidth}\\[1mm]
\linead{\linewidth}

\begin{infoband}{milcNegro}
\textbf{Nota para el docente.} Las tres son citas textuales y llevan su referencia debajo. Puedes remitir a la obra en clase.
\end{infoband}

\milcestacion{milcVino}{Mi autoevaluación · marca una casilla por fila}

{\small
\setlength{\extrarowheight}{2.2pt}
\begin{tabularx}{\textwidth}{>{\RaggedRight\arraybackslash\bfseries}p{3.0cm}YYY}
\rowcolor{milcVino}\color{white}\bfseries Criterio & \color{white}\bfseries Lo logré & \color{white}\bfseries En proceso & \color{white}\bfseries Necesito apoyo\\[.6mm]
Comprensión del tema & \milcopcion{Explico las ideas con mis palabras.} & \milcopcion{Reconozco las ideas, falta precisión.} & \milcopcion{Necesito repasar lo básico.}\\[.8mm]
\rowcolor{milcGris}Calidad del reparto & \milcopcion{Separo lo que depende de mí de lo que no, y uso una estrategia distinta para cada columna.} & \milcopcion{Hago el reparto, pero sigo intentando controlar lo de la segunda columna.} & \milcopcion{Sigo pensando que preocuparse mucho es una forma de resolver.}\\[.8mm]
Aplicación y producto & \milcopcion{Mi producto cumple los criterios.} & \milcopcion{Está armado, con ajustes pendientes.} & \milcopcion{Aún está en construcción.}\\[.8mm]
\rowcolor{milcGris}Reflexión 5 dimensiones & \milcopcion{Respondí cada una con honestidad.} & \milcopcion{Respondí algunas con detalle.} & \milcopcion{Mis respuestas son muy generales.}\\[.8mm]
Triángulo de pensamiento & \milcopcion{Conecté las 3 voces con mi trabajo.} & \milcopcion{Conecté al menos una.} & \milcopcion{No las relaciono con mi proceso.}\\[.8mm]
\end{tabularx}}

\vspace{3mm}
\textbf{Hoy entendí que:}\\[1mm]
\linead{\linewidth}\\[1mm]
\linead{\linewidth}

\textbf{Mi compromiso para la próxima sesión es:}\\[1mm]
\linead{\linewidth}\\[1mm]
\linead{\linewidth}

\begin{softbox}{milcMostaza}{milcGris}{Cita de despedida · Marco Aurelio}
\textit{``El obstáculo en el camino se vuelve el camino. Lo que estorba al camino se vuelve camino.''}\\[1mm]
Cada error de hoy, bien observado, se convierte en pieza del camino que pisarás mañana.
\end{softbox}

\small
\begin{tabularx}{\textwidth}{>{\bfseries}p{3.6cm}Y}
\rowcolor{milcGradeDark!12}Nombre completo: & \linead{10cm}\\
Fecha: & \linead{4cm} \quad Curso/grupo: \linead{4cm}\\
Mi firma: & \linead{9cm}\\
\end{tabularx}
\normalsize

\begin{infoband}{milcGradeDark}
\textbf{Para la próxima sesión.} Dos cosas. \textbf{Primero, haz \emph{una} acción de la columna izquierda} ---la que tenga la fecha más cercana--- y anota qué pasó con la preocupación después de hacerla. \emph{Casi siempre baja, y no porque el problema se resuelva: baja porque dejó de estar pendiente.} \textbf{Segundo, pregunta en tu casa por la roya o por otra que se les haya venido encima:} averigua con alguien mayor \textbf{una época en que llegó algo que nadie podía evitar} ---una plaga, una creciente, una quiebra, una enfermedad, un desplazamiento--- y pregúntale \textbf{qué se pudo hacer y qué no}, y \textbf{qué decisión sí estuvo en sus manos}. \textbf{Trae:} el resultado de tu acción y lo que te contaron. \emph{Vas a encontrar la tabla de esta guía contada en primera persona: casi todos distinguen con mucha claridad lo que les cayó encima de lo que decidieron, y casi todos recuerdan mejor la decisión que la desgracia.}
\end{infoband}

\end{document}
